\documentclass[12pt, fleqn]{article}

\usepackage[utf8]{inputenc}
\usepackage[T2A]{fontenc}
\usepackage{amssymb, amsmath, mathrsfs, amsthm}
\usepackage[russian]{babel}
\usepackage{graphicx}
\usepackage[footnotesize]{caption2}
\usepackage{indentfirst}
\usepackage[hidelinks]{hyperref}
\usepackage{multirow}
\usepackage{pgfplotstable}
\usepackage{makecell}
%\usepackage[ruled,section]{algorithm}
%\usepackage[noend]{algorithmic}
%\usepackage[all]{xy}

% Параметры страницы
\textheight=24cm
\textwidth=16cm
\oddsidemargin=5mm
\evensidemargin=-5mm
\marginparwidth=36pt
\topmargin=-1cm
\footnotesep=3ex
%\flushbottom
\raggedbottom
\tolerance 3000
% подавить эффект "висячих стpок"
\clubpenalty=10000
\widowpenalty=10000
\renewcommand{\baselinestretch}{1.1}
\renewcommand{\baselinestretch}{1.5} %для печати с большим интервалом

% Дополнительные команды для личных обозначений
\newcommand{\expectation}{\mathop{\mathbb{E}}}
\newcommand{\norm}[1]{\left\lVert#1\right\rVert}
\newcommand{\loss}{\mathop{\mathcal{L}}}
\newcommand{\mse}{\mathop{MSE}}
\newcommand{\scalarproduct}[1]{\langle #1 \rangle}

\newcommand{\predictionfunction}{\hat{f}}
\newcommand{\ensemblefunction}{a}
\newcommand{\optimizationmethodfunction}{\tilde{f}}
\newcommand{\distinguishparameter}{\alpha}
\newcommand{\objects}{X}
\newcommand{\results}{Y}
\newcommand{\predictedobjects}{\widehat{\objects}}

\newcommand{\numberobjects}{N}
\newcommand{\numberpredictionfunctions}{M}

\newcommand{\for}[3]{\sum\limits_{#1 = #2}^{#3}}  % Usage: \for{index}{begin}{end}
\newcommand{\forn}[2]{\for{#1}{1}{#2}}  % Usage: \forn{index}{end}

\newcommand{\many}[3]{#1 1 #2, #1 2 #2, \dots, #1 #3 #2}  % Usage: \many{prefix}{suffix}{end}

\newcommand{\reference}[1]{(\hyperref[#1]{\ref{#1}})}

\newcommand{\ensemblefunctionfull}{\ensemblefunction(\many{\predictionfunction_}{(x)}{\numberpredictionfunctions})}

\begin{document}

\begin{titlepage}
\begin{center}
    Московский государственный университет имени М. В. Ломоносова

    \bigskip
    \includegraphics[width=50mm]{msu.eps}

    \bigskip
    Факультет Вычислительной Математики и Кибернетики\\
    Кафедра Математических Методов Прогнозирования\\[10mm]

    \textsf{\large\bfseries
        <<Сравнение скорости вычисления собственных значений\\ положительно определённых матриц\\ при помощи QR алгоритма>>
    }\\[20mm]

    \begin{flushright}
        \parbox{0.5\textwidth}{
            Выполнил:\\
            студент 1 курса магистратуры 517 группы\\
            \emph{Королев Николай Сергеевич}\\[5mm]
            Преподаватель:\\
            канд. техн. наук, доцент\\
            \emph{Русол Андрей Владимирович}
        }
    \end{flushright}

    \vspace{\fill}
    Москва, 2021
\end{center}
\end{titlepage}

\newpage
\section{Постановка задачи}

Исследовать способы ускорения вычисления собственных значений положительно определённых матриц.

\subsection{Постановка задачи о вычислении собственных значений положительно определённой матрицы}
Дана положительно определённая матрица $A$ размера $N \times N$. Необходимо вычислить все $N$ её собственных значений $\lambda_1, \lambda_2, \dots, \lambda_N$.

\section{QR алгоритм}

Для нахождения всех собственных значений положительно определённой матрицы $A$ можно воспользоваться QR алгоритмом, который выглядит следующим образом:
\begin{enumerate}
	\item Обозначим $A_0 := A, \; k := 0$.
	\item Представить матрицу $A_k$ в виде произведения унитарной матрицы $Q_k$ и верхнетреугольной матрицы $R_k$. $Q_k R_k = A_k$
	\item Вычислить $A_{k+1} := R_k Q_k$
	\item Увеличить $k$ на единицу. $k := k + 1$
	\item Повторить шаги 2-4 до тех пор пока внедиагональные элементы матрицы $A_k$ не станут близкими к нулю.
	\item Значения на диагонали матрицы $A_k$ будут являться приближением собственных значений матрицы $A$.
\end{enumerate}

\subsection{Доказательство корректности алгоритма}
Заметим, что все матрицы $A_k$ для $k = 0, 1, \dots$ являются подобными, т.к. \\
$A_{k + 1} = R_k Q_k = Q_k^{-1} Q_k R_k Q_k = Q_k^{-1} A_k Q_k = Q_k^T A_k Q_k$, а значит их собственные значения совпадают.

Также для положительно определённой матрицы $A$ известно \cite{QR}, что внедиагольные элементы матрицы $A_k$ будут стремиться к нулю при $k \to \infty$.

\section{Вычислительные эксперименты}

QR алгоритм был реализован тремя различными способами на языке Python 3 при помощи библиотеки Numpy для использования векторизации вычислений, после чего лучшая из имплементаций была ускорена при помощи JIT-компилятора Numba. Также алгоритм был имплементирован на языке программирования C++, но без использования параллельных вычислений. Результаты измерений для различных размеров матрицы $A$ приведены в таблице \ref{tabular:results}. Все имплементации выложены в открытый доступ на Github\footnote{\url{https://github.com/CrafterKolyan/eigenvalues-speed-comparison/blob/main/experiments/}}.

\begin{center}
	\pgfplotstabletypeset[
	after row=\hline,
	col sep=semicolon,
	string type,
	columns/index/.style={column name={Имплементация
		}, column type={|c}},
	columns/10/.style={column name={$N = 10$}, column type={|c}},
	columns/20/.style={column name={$N = 20$}, column type={|c}},
	columns/30/.style={column name={$N = 30$}, column type={|c}},
	columns/40/.style={column name={$N = 40$}, column type={|c}},
	columns/50/.style={column name={$N = 50$}, column type={|c|}},
	every head row/.style={before row=\hline,after row=\hline},
	every last row/.style={after row=\hline},
	]{tables/results.csv}
	
	\begin{table}[h!]
		\caption{Время выполнения (в миллисекундах) QR алгоритма различных имплементаций при различных размерах исходной матрицы. В таблице приведено среднее время выполнения ± среднеквадратичное отклонение времени по 7 замерам времени, в каждом из которых алгоритм запускался 100 раз.}
		\label{tabular:results}
	\end{table}
\end{center}

\subsection{Анализ полученных результатов}

Полученные результаты показывают, что JIT-компилятор Numba позволяет добиться ускорения примерно в 4 раза в сравнении с обычным использованием Numpy, выигрывая по скорости на больших матрицах даже у C++. Что не менее важно, для того, чтобы воспользоваться JIT-компилятором, не требуется тратить дополнительное время на разработку решения. Также видно, что при этом компилятор Numba даже не нужно дополнительно настраивать, т.к. в любом случае получается выигрыш в скорости работы.

Кроме того, отметим, что библиотечная функция из Numpy (\texttt{numpy.linalg.eigvals}) работает примерно в 10 раз быстрее, чем самая быстрая из реализаций представленных в работе, но неизвестно какой конкретно алгоритм применяется в библиотечной функции, поэтому сравнение скоростей работы функции из библиотеки и представленных имплементаций может быть некорректно.

\section{Заключение}

В процессе выполнения работы были получены следующие результаты:
\begin{itemize}
\item Использование JIT-компилятора Numba позволяет получить повышение скорости работы в 4 раза по сравнению с обычным Numpy при реализации QR~алгоритма;
\item JIT-компилятор Numba вместе с использованием Numpy позволяет получить более быструю реализацию QR алгоритма на больших матрицах, нежели чем аналогичная реализация на C++ без параллелизации, кроме того, время затраченное на написание кода с использованием Numba и Numpy сильно меньше, чем реализация аналогичного решения на C++.
\end{itemize}

\def\BibUrl#1.{}\def\BibAnnote#1.{}
%\def\BibUrl#1{\\{\footnotesize\tt\def~{\char126} http://#1}}
\bibliographystyle{gost71s}
\bibliography{bibliography}

\end{document}
